\documentclass[12pt]{article}
\usepackage{anysize}
\usepackage{xltxtra}
\usepackage{xunicode}
\usepackage{graphicx}
\usepackage{color}
\usepackage{tikz}
\usepackage{minted}
\usepackage{listings}
\usepackage{enumitem}
\usepackage{framed}
\usepackage{relsize}
\usepackage{float}
\usepackage{algpseudocode}
\usepackage[hyphens]{url}

\usepackage{color,hyperref}

\usepackage{fancyhdr}
\usepackage{lastpage}
\usepackage{fontspec}

\linespread{1.5}

\definecolor{darkblue}{rgb}{0.0,0.0,0.3}
\hypersetup{colorlinks,breaklinks,
        linkcolor=darkblue,urlcolor=darkblue,
            anchorcolor=black,citecolor=black}

\def\name{Grigorios Lyras}
\def\IF{\textbf{if}}
\def\THEN{\textbf{then}}
\def\SAYS{\textbf{says}}
\def\SPEAKSFOR{$\Rightarrow$}
\def\FOR{\textbf{for}}
\def\AND{\textbf{and}}

\setlength{\headheight}{15.2pt}

\pagestyle{fancy}{ %
    \fancyhf{} % remove everything
    \renewcommand{\headrulewidth}{0pt} % remove lines as well
    \renewcommand{\footrulewidth}{0.5pt}
    \lfoot{\name}
    \rfoot{Page \thepage\ of \pageref{LastPage}}

\renewcommand{\labelenumi}{(\alph{enumi})}
\defaultfontfeatures{Ligatures=TeX} % To support LaTeX quoting style
\setmainfont[Mapping=TeX-text]{CMU Serif}


\begin{document}

\begin{titlepage}

\begin{tikzpicture}[remember picture, overlay]
      \node [anchor=north east, inner sep=0pt]  at (current page.north east)
           {\includegraphics[height=3cm]{title/ucl_logo}};
       \end{tikzpicture}
\begin{center}
\begin{LARGE}\textbf{
University College London\\}\end{LARGE}
\vspace{1cm}
\begin{Large}Simulating Path Planning for Data\end{Large}\\
\begin{Large}Muling in Wireless Sensor Networks\end{Large}\\
\begin{Large}using Unmanned Aerial Vehicles\end{Large}\\
\vspace{2cm}
\begin{LARGE}
MSc Networked Computer Systems\\
NCS Group Project Individual Report\\
\end{LARGE}
\large{Academic Year 2013-2014}\\
\vfill
\end{center}
This report is submitted as part requirement for the MSc in Networked Computer
Systems at University College London. It is substantially the result of our own
work except where explicitly indicated in the text. The report may be freely
copied and distributed provided the source is explicitly acknowledged.\\
\begin{center}
\begin{tabular}{l r}
\Large{Grigorios Lyras}&
\large{SID: 13047325}\\
\end{tabular}\\
\vspace{1cm}
\large\today\\
\end{center}
\end{titlepage}



\tableofcontents{}

\pagebreak

\section{Introduction}
The aim of this report is the discussion of the Group Project that I took part
in as part of the NCS team, comprised of myself and Seng Jea Lee, Johann Mifsud,
Minmin Zhang. In this report I will present my own role in the progress of
the project. The analysis will be broken into three main sections. In objectives
and achievements I will present the objectives we set through the duration of
the project as well as the achievements and my contribution to them. Moving to
the results of the project we will discuss the output of the project as well as
the quality of the outcome. In the third section, I will present my view on the
project management techniques and tools that we used along with my experience on
them. The interaction with the rest of the team members and the task
distribution will also be reported. In the final section I will present my own
contribution to the project.

\section{Project Overview}
Wireless Sensor Networks are systems of sensors distributed over an area in
order to collect useful information. These are low cost devices with limitations
in their capabilities. One of these limitations is the communication range. As a
result it is understood that these networks are not always connected. A concept
to overcome such difficulties is the use of data mules that visit these devices
and collect the data from them. These data mules are networked devices that are
able to move from location to location, pick up the data, store them
temporarily and then deliver them back to the base. An interesting area is the
exploration of data muling with the use of Unmanned Aerial Vehicles (UAVs).
UAVs have unique mobility abilities which make them a very suitable solution
for data muling.

Since the suitability of UAVs for the task is established we move to the
implementation. In order to ensure that UAVs behave in a safe manner thorough
testing is required. However this is not always possible due to weather
conditions. Another factor that increases the risk significantly in this
research area is the cost of the UAVs the nearby environment. The high cost
makes field testing unsuitable for this research area. As a consequence the
research has focused on simulating UAVs in order to reduce the risk mentioned
above.

The initial idea for this project was the exploration of path planning
algorithms for UAVs when they are used to collect data from distributed wireless
sensor networks. Due to the limitations and risk involved in such a project if it
is taken on the field the exploration was limited to a simulated environment.
After researching the area we focused on a simulator that enables the physical
simulation of UAVs. CRATES was the simulator we used for this project which
lacked support for modeling the wireless channel.

We identified the areas where the project would contribute to. This project set
to build a radio propagation model for CRATES and utilise this in order to
explore higher level Path Planning Algorithms. These were to be evaluated on
their performance based on distance and data collection statistics.


\section{Objectives and Achievements}
One of the most challenging aspects of this project has been it's wide scope. It
is in effect a multi disciplinary project. The main objectives were the
extension of CRATES in order to support wireless communication as well as
the implementation and testing of path planning algorithms. However, these two
main objectives were divided into multiple milestones along the progress of the
project.

\subsection{Wireless Channel Modeling}
Taking into account the project outline one of the first milestones along the
timeline was the Modeling of the wireless channel in CRATES. This task was split
into two main subgoals. Firstly we had to comprehend how CRATES operated and how
it would be best to implement the desired functionality. In parallel we were to
research on how to model the wireless channel.

CRATES supports wireless sensors however, it is limited because it offers no
communication model. In order to add support for the communication model we
extended CRATES and enabled the models inside the simulated world to pass
messages from one to another. These communiation primitives were to be enriched
with a realistic radio model that would control the radio propagation method.
Since we required a detailed wireless communication simulation, it was
essential to simulate not only the raw wireless channel, but also the network
stack layers up to the internet protocol.

In order to simulate the network stack layers we reused existing code. In
particular we used ns2, a disrete event simulator, that executes TCL scripts at
runtime. ns2 is a widely used simulator in researh areas and supports multiple
MAC layers. Alternatives to ns2 were considered, however, due to our previous
familiarity with ns2 was what led us to use ns2 for our project. In our
implementation we run two simulators in parallel. We treat the ns2 simulation
as a funtion call, pushing the input, waiting for a fixed amount of time and
then parsing the result.

In the area of wireless modeling we had to gather data for the network modules
that were available on the real life platforms. To do that we created a
measurement platform using OrisenStamps and NVO8C GPS units that allowed us to
measure wireless data as well as GPS coordinates. This way we could calculate
the distanes between the different platforms.

\subsection{Path Planning}
The second main goal for our project has been the exploration of path planning
algorithms. We focused on two main algorithms; one based on Sugihara's path
planning algorithm and a modified genetic algorithm. These two algorithms are
tested and their performance is compared on the basis of total distance
travelled by the UAVs and the amount of data colleted.

\section{Project Results}
In this section of the report
\subsection{Radio Modeling}
\subsubsection{ZigBee}
\subsubsection{WiFi}
\subsection{Simulator Plugin}
\subsection{Path Planning}
a critical assessment of the project results (e.g., how do you judge the
quality of the final result, what were the important technical decisions
this project, what worked well, what decisions you would change in the light
of experience, etc.);

\section{Project Management}
a critical assessment of the project management techniques used (e.g.,
whether you followed standard software engineering processes, and if so then
how well they worked or didn't work);

a discussion and critical assessment of how well the team and each of its
members worked from your point-of-view (e.g., whether decisions were
adequately agreed, specified and documented, or had to be revised later,
whether the work was appropriately split among team members);

a description and critical assessment of your own workplan and contribution
to the project (e.g., what aspects of the project did you contribute most
to; whether you were over-worked or under-worked; whether you as an
individual could have achieved more with different approaches and/or
organisation).

In the Individual report, you may make reference to the submitted Group
documents for further detail, but the report should stand on its own.

Note that the Individual report is only made available to the Examiners and will
not be seen by any other students.


\end{document}
